\sectionkicker{Source-backed technical notes}
\subsection*{Half-year Synthesis — 進歩の単位はModelからExecution Stackへ: Technical Notes}
本欄は記事本文で圧縮した一次資料上の情報を、比較・再検証しやすい形で整理したものである。時系列、確認済みの主張、留意点、一次資料URLを掲載する。
\medskip
\subsection*{Theme at a glance}
\addcontentsline{toc}{subsection}{Theme at a glance}
\begin{center}
\footnotesize
\begin{tabularx}{\linewidth}{@{}p{0.20\linewidth}p{0.10\linewidth}p{0.14\linewidth}X@{}}
\toprule
資料 & 位置づけ & 種別 & 時系列 \\
\midrule
ChatGPT agent & 補足資料 & Agent & 2025-07-17 (Agent（公開）); 2025-07-17 (System Card公開)  \\
Claude 4.5 model and agent platform sequence & 補足資料 & モデル更新 & 2025-09-29 (CLAUDE SONNET 4 5（公開）); 2025-10-15 (CLAUDE HAIKU 4 5（公開）); 2025-11-24 (CLAUDE OPUS 4 5（公開）)  \\
Gemini API H2 2025 lifecycle & 補足資料 & API & 2025-07-07 (BATCH MODE（公開）); 2025-07-17 (VEO3（Preview）); 2025-07-22 (GEMINI 2 5 FLASH LITE（公開）); 2025-08-14 (IMAGEN4（一般提供）); 2025-08-26 (GEMINI 2 5 IMAGE（Preview）); 2025-09-09 (VEO3（一般提供）); 2025-09-25 (ROBOTICS ER 1 5（Preview）); 2025-10-07 (COMPUTER USE（Preview）); 2025-10-15 (VEO3 1（Preview）); 2025-11-18 (GEMINI3 PRO（Preview）); 2025-12-11 (INTERACTIONS API AND DEEP 研究Preview（技術イベント）); 2025-12-17 (GEMINI3 FLASH（Preview）)  \\
GPT-5 release / developer API / safe-completions & 補足資料 & モデル & 2025-08-07 (モデル公開); 2025-08-07 (API公開); 2025-08-07 (SAFETY METHOD（公開）)  \\
gpt-oss-120b / gpt-oss-20b & 補足資料 & オープンウェイト & 2025-08-05 (オープンウェイト公開); 2025-08-05 (モデル CARD（公開）)  \\
Google DeepMind H2 world-model, robotics, and agent research & 補足資料 & モデル更新 & 2025-08-05 (GENIE 3 ANNOUNCEMENT（技術イベント）); 2025-09-25 (GEMINI ROBOTICS 1 5 ANNOUNCEMENT（技術イベント）); 2025-11-13 (SIMA 2 ANNOUNCEMENT（技術イベント）); 2025-11-18 (GEMINI 3 ANNOUNCEMENT（技術イベント）)  \\
vLLM v0.10→v0.13 & 補足資料 & Framework & 2025-07-24 (RUNTIME（公開）); 2025-10-02 (RUNTIME（公開）); 2025-12-19 (RUNTIME（公開）)  \\
OpenAI–Anthropic joint safety evaluation & 補足資料 & 安全性関連 & 2025-08-27 (CROSS LAB SAFETY EVALUATION（技術イベント）)  \\
\bottomrule
\end{tabularx}
\end{center}
\normalsize

\begin{technicalnote}{ChatGPT agent}{補足資料}
\begingroup
% reader-facing Technical Notes break policy
\widowpenalty=10000
\clubpenalty=10000
\displaywidowpenalty=10000
\begin{tabularx}{\linewidth}{@{}>{\bfseries}p{0.22\linewidth}X@{}}
組織 & OpenAI \\
種別 & Agent
時系列 & 2025-07-17 (Agent（公開）); 2025-07-17 (System Card公開) \\
\end{tabularx}
\smallskip
{\bfseries 一次資料から整理したtechnical points}
\begin{itemize}[leftmargin=1.5em,itemsep=0.35em]
\item \textbf{一次情報で確認できる事実}: 「ChatGPT agent」について、一次資料で確認できる公開・提供・機能・時系列上の事実を要約した項目。
\item \textbf{一次情報で確認できる事実}: 「ChatGPT agent」について、一次資料で確認できる公開・提供・機能・時系列上の事実を要約した項目。
\end{itemize}
\begin{minipage}{\linewidth}
% reader-facing Technical Notes generic boundary/source tail group
{\bfseries 読む際の境界}
\begin{itemize}[leftmargin=1.5em,itemsep=0.35em]
\item \textbf{分析上の留意点}: 「ChatGPT agent」の扱いでは、一次資料で確認できる範囲、提供時点、評価条件を維持し、記載範囲を超えて一般化しない。
\end{itemize}
\begin{samepage}
{\bfseries 一次資料}
\begingroup\sloppy
\Urlmuskip=0mu plus 2mu\relax
\begin{itemize}[leftmargin=1.5em,itemsep=0.25em]
\item {\scriptsize\url{https://openai.com/index/chatgpt-agent-system-card}}
\item {\scriptsize\url{https://openai.com/index/introducing-chatgpt-agent}}
\end{itemize}
\endgroup
\end{samepage}
\end{minipage}
\endgroup
\end{technicalnote}
\medskip

\begin{technicalnote}{Claude 4.5 model and agent platform sequence}{補足資料}
\begingroup
% reader-facing Technical Notes break policy
\widowpenalty=10000
\clubpenalty=10000
\displaywidowpenalty=10000
\begin{tabularx}{\linewidth}{@{}>{\bfseries}p{0.22\linewidth}X@{}}
組織 & Anthropic \\
種別 & モデル更新
時系列 & 2025-09-29 (CLAUDE SONNET 4 5（公開）); 2025-10-15 (CLAUDE HAIKU 4 5（公開）); 2025-11-24 (CLAUDE OPUS 4 5（公開）) \\
\end{tabularx}
\smallskip
{\bfseries 一次資料から整理したtechnical points}
\begin{itemize}[leftmargin=1.5em,itemsep=0.35em]
\item \textbf{一次情報で確認できる事実}: Anthropicは9月29日、Claude Sonnet 4.5を、Claude Codeの更新およびClaude Agent SDKとあわせて公開した。
\item \textbf{一次情報で確認できる事実}: Anthropicは10月15日にClaude Haiku 4.5、11月24日にClaude Opus 4.5を公開し、いずれもAPIで提供した。
\item \textbf{分析上の留意点}: Anthropicの2025年後半の系列では、モデル公開がコーディング、Computer Use、コンテキスト・メモリ、エージェント構築の各利用面と結び付いて展開された。
\end{itemize}
\begin{minipage}{\linewidth}
% reader-facing Technical Notes generic boundary/source tail group
{\bfseries 読む際の境界}
\begin{itemize}[leftmargin=1.5em,itemsep=0.35em]
\item \textbf{分析上の留意点}: 保存したnewsroomのスナップショットはページ分割された継続更新型indexで、先頭HTMLは主に2026年の項目を示していた。2025年後半の年表は、個別のAnthropic公式ページを追跡して再構成した。
\end{itemize}
\begin{samepage}
{\bfseries 一次資料}
\begingroup\sloppy
\Urlmuskip=0mu plus 2mu\relax
\begin{itemize}[leftmargin=1.5em,itemsep=0.25em]
\item {\scriptsize\url{https://www.anthropic.com/news}}
\end{itemize}
\endgroup
\end{samepage}
\end{minipage}
\endgroup
\end{technicalnote}
\medskip

\begin{technicalnote}{Gemini API H2 2025 lifecycle}{補足資料}
\begingroup
% reader-facing Technical Notes break policy
\widowpenalty=10000
\clubpenalty=10000
\displaywidowpenalty=10000
\begin{tabularx}{\linewidth}{@{}>{\bfseries}p{0.22\linewidth}X@{}}
組織 & Google \\
種別 & API
時系列 & 2025-07-07 (BATCH MODE（公開）); 2025-07-17 (VEO3（Preview）); 2025-07-22 (GEMINI 2 5 FLASH LITE（公開）); 2025-08-14 (IMAGEN4（一般提供）); 2025-08-26 (GEMINI 2 5 IMAGE（Preview）); 2025-09-09 (VEO3（一般提供）); 2025-09-25 (ROBOTICS ER 1 5（Preview）); 2025-10-07 (COMPUTER USE（Preview）); 2025-10-15 (VEO3 1（Preview）); 2025-11-18 (GEMINI3 PRO（Preview）); 2025-12-11 (INTERACTIONS API AND DEEP 研究Preview（技術イベント）); 2025-12-17 (GEMINI3 FLASH（Preview）) \\
\end{tabularx}
\smallskip
{\bfseries 一次資料から整理したtechnical points}
\begin{itemize}[leftmargin=1.5em,itemsep=0.35em]
\item \textbf{一次情報で確認できる事実}: Gemini APIの変更履歴には、非同期バッチ処理、Embedding、動画、画像、ネイティブ音声、ロボティクス、Computer Use、Gemini 3、統合Interactions API、Deep Research agentにまたがる高密度な2025年後半の更新が記録されている。
\item \textbf{一次情報で確認できる事実}: 変更履歴はPreview、一般提供、非推奨化、停止を明示的に区別しており、これらの提供状態自体が技術年表の一部である。
\item \textbf{一次情報で確認できる事実}: Gemini 3 Pro Previewは11月18日にAPIへ追加され、Gemini 3 Flash Previewは12月17日に続いた。
\end{itemize}
\begin{minipage}{\linewidth}
% reader-facing Technical Notes generic boundary/source tail group
{\bfseries 読む際の境界}
\begin{itemize}[leftmargin=1.5em,itemsep=0.35em]
\item \textbf{分析上の留意点}: この変更履歴はサービス／APIの年表であり、上流の研究・モデル公開年表ではない。Previewから一般提供への移行や停止を、最初のモデル発表と同一イベントへ統合しない。
\end{itemize}
\begin{samepage}
{\bfseries 一次資料}
\begingroup\sloppy
\Urlmuskip=0mu plus 2mu\relax
\begin{itemize}[leftmargin=1.5em,itemsep=0.25em]
\item {\scriptsize\url{https://ai.google.dev/gemini-api/docs/changelog}}
\end{itemize}
\endgroup
\end{samepage}
\end{minipage}
\endgroup
\end{technicalnote}
\medskip

\begin{technicalnote}{GPT-5 release / developer API / safe-completions}{補足資料}
\begingroup
% reader-facing Technical Notes break policy
\widowpenalty=10000
\clubpenalty=10000
\displaywidowpenalty=10000
\begin{tabularx}{\linewidth}{@{}>{\bfseries}p{0.22\linewidth}X@{}}
組織 & OpenAI \\
種別 & モデル
時系列 & 2025-08-07 (モデル公開); 2025-08-07 (API公開); 2025-08-07 (SAFETY METHOD（公開）) \\
\end{tabularx}
\smallskip
{\bfseries 一次資料から整理したtechnical points}
\begin{itemize}[leftmargin=1.5em,itemsep=0.35em]
\item \textbf{一次情報で確認できる事実}: 「GPT-5 release / developer API / safe-completions」について、一次資料で確認できる公開・提供・機能・時系列上の事実を要約した項目。
\item \textbf{一次情報で確認できる事実}: 「GPT-5 release / developer API / safe-completions」について、一次資料で確認できる公開・提供・機能・時系列上の事実を要約した項目。
\end{itemize}
\begin{minipage}{\linewidth}
% reader-facing Technical Notes generic boundary/source tail group
{\bfseries 読む際の境界}
\begin{itemize}[leftmargin=1.5em,itemsep=0.35em]
\item \textbf{分析上の留意点}: 「GPT-5 release / developer API / safe-completions」の扱いでは、一次資料で確認できる範囲、提供時点、評価条件を維持し、記載範囲を超えて一般化しない。
\end{itemize}
\begin{samepage}
{\bfseries 一次資料}
\begingroup\sloppy
\Urlmuskip=0mu plus 2mu\relax
\begin{itemize}[leftmargin=1.5em,itemsep=0.25em]
\item {\scriptsize\url{https://openai.com/index/gpt-5-safe-completions}}
\item {\scriptsize\url{https://openai.com/index/introducing-gpt-5}}
\item {\scriptsize\url{https://openai.com/index/introducing-gpt-5-for-developers}}
\end{itemize}
\endgroup
\end{samepage}
\end{minipage}
\endgroup
\end{technicalnote}
\medskip

\begin{technicalnote}{gpt-oss-120b / gpt-oss-20b}{補足資料}
\begingroup
% reader-facing Technical Notes break policy
\widowpenalty=10000
\clubpenalty=10000
\displaywidowpenalty=10000
\begin{tabularx}{\linewidth}{@{}>{\bfseries}p{0.22\linewidth}X@{}}
組織 & OpenAI \\
種別 & オープンウェイト
時系列 & 2025-08-05 (オープンウェイト公開); 2025-08-05 (モデル CARD（公開）) \\
\end{tabularx}
\smallskip
{\bfseries 一次資料から整理したtechnical points}
\begin{itemize}[leftmargin=1.5em,itemsep=0.35em]
\item \textbf{一次情報で確認できる事実}: 「gpt-oss-120b / gpt-oss-20b」について、一次資料で確認できる公開・提供・機能・時系列上の事実を要約した項目。
\item \textbf{Vendor claim}: 「gpt-oss-120b / gpt-oss-20b」について、提供元・プロジェクト・著者側の評価または説明として記録された項目。独立再現された結果を意味しない。
\end{itemize}
\begin{minipage}{\linewidth}
% reader-facing Technical Notes generic boundary/source tail group
{\bfseries 読む際の境界}
\begin{itemize}[leftmargin=1.5em,itemsep=0.35em]
\item \textbf{分析上の留意点}: 「gpt-oss-120b / gpt-oss-20b」の扱いでは、一次資料で確認できる範囲、提供時点、評価条件を維持し、記載範囲を超えて一般化しない。
\end{itemize}
\begin{samepage}
{\bfseries 一次資料}
\begingroup\sloppy
\Urlmuskip=0mu plus 2mu\relax
\begin{itemize}[leftmargin=1.5em,itemsep=0.25em]
\item {\scriptsize\url{https://openai.com/index/gpt-oss-model-card}}
\item {\scriptsize\url{https://openai.com/index/introducing-gpt-oss}}
\end{itemize}
\endgroup
\end{samepage}
\end{minipage}
\endgroup
\end{technicalnote}
\medskip

\begin{technicalnote}{Google DeepMind H2 world-model, robotics, and agent research}{補足資料}
\begingroup
% reader-facing Technical Notes break policy
\widowpenalty=10000
\clubpenalty=10000
\displaywidowpenalty=10000
\begin{tabularx}{\linewidth}{@{}>{\bfseries}p{0.22\linewidth}X@{}}
組織 & Google DeepMind \\
種別 & モデル更新
時系列 & 2025-08-05 (GENIE 3 ANNOUNCEMENT（技術イベント）); 2025-09-25 (GEMINI ROBOTICS 1 5 ANNOUNCEMENT（技術イベント）); 2025-11-13 (SIMA 2 ANNOUNCEMENT（技術イベント）); 2025-11-18 (GEMINI 3 ANNOUNCEMENT（技術イベント）) \\
\end{tabularx}
\smallskip
{\bfseries 一次資料から整理したtechnical points}
\begin{itemize}[leftmargin=1.5em,itemsep=0.35em]
\item \textbf{一次情報で確認できる事実}: Google DeepMindは8月5日、リアルタイムで対話可能なworld modelとしてGenie 3を発表した。
\item \textbf{一次情報で確認できる事実}: Google DeepMindは9月25日、物理行動と身体化推論・ツール利用を組み合わせるGemini Robotics 1.5とRobotics-ER 1.5を発表した。Robotics-ER 1.5はGemini APIで提供され、Robotics 1.5は一部パートナー向けに限定された。
\item \textbf{一次情報で確認できる事実}: 11月13日のSIMA 2はGeminiの推論を3D世界のエージェントへ統合し、11月18日にはより広範なfrontier modelとしてGemini 3が続いた。
\end{itemize}
\begin{minipage}{\linewidth}
% reader-facing Technical Notes generic boundary/source tail group
{\bfseries 読む際の境界}
\begin{itemize}[leftmargin=1.5em,itemsep=0.35em]
\item \textbf{分析上の留意点}: 保存したDeepMindのindexスナップショットはページ分割されており、先頭ページだけでは2025年後半の完全な記録にならないため、個別の公式ページを用いて検証した。
\item \textbf{分析上の留意点}: 研究プロトタイプ、一部パートナー限定のロボティクス提供、一般提供されるAPIモデルは、それぞれ異なる提供状態として区別する。
\end{itemize}
\begin{samepage}
{\bfseries 一次資料}
\begingroup\sloppy
\Urlmuskip=0mu plus 2mu\relax
\begin{itemize}[leftmargin=1.5em,itemsep=0.25em]
\item {\scriptsize\url{https://deepmind.google/blog/}}
\end{itemize}
\endgroup
\end{samepage}
\end{minipage}
\endgroup
\end{technicalnote}
\medskip

\begin{technicalnote}{vLLM v0.10→v0.13}{補足資料}
\begingroup
% reader-facing Technical Notes break policy
\widowpenalty=10000
\clubpenalty=10000
\displaywidowpenalty=10000
\begin{tabularx}{\linewidth}{@{}>{\bfseries}p{0.22\linewidth}X@{}}
組織 & vLLM project \\
種別 & Framework
時系列 & 2025-07-24 (RUNTIME（公開）); 2025-10-02 (RUNTIME（公開）); 2025-12-19 (RUNTIME（公開）) \\
\end{tabularx}
\smallskip
{\bfseries 一次資料から整理したtechnical points}
\begin{itemize}[leftmargin=1.5em,itemsep=0.35em]
\item \textbf{Project claim}: 「vLLM v0.10→v0.13」について、提供元・プロジェクト・著者側の評価または説明として記録された項目。独立再現された結果を意味しない。
\end{itemize}
\begin{minipage}{\linewidth}
% reader-facing Technical Notes generic boundary/source tail group
{\bfseries 読む際の境界}
\begin{itemize}[leftmargin=1.5em,itemsep=0.35em]
\item \textbf{分析上の留意点}: 「vLLM v0.10→v0.13」の扱いでは、一次資料で確認できる範囲、提供時点、評価条件を維持し、記載範囲を超えて一般化しない。
\end{itemize}
\begin{samepage}
{\bfseries 一次資料}
\begingroup\sloppy
\Urlmuskip=0mu plus 2mu\relax
\begin{itemize}[leftmargin=1.5em,itemsep=0.25em]
\item {\scriptsize\url{https://github.com/vllm-project/vllm/releases/tag/v0.10.0}}
\item {\scriptsize\url{https://github.com/vllm-project/vllm/releases/tag/v0.11.0}}
\item {\scriptsize\url{https://github.com/vllm-project/vllm/releases/tag/v0.13.0}}
\end{itemize}
\endgroup
\end{samepage}
\end{minipage}
\endgroup
\end{technicalnote}
\medskip

\begin{technicalnote}{OpenAI–Anthropic joint safety evaluation}{補足資料}
\begingroup
% reader-facing Technical Notes break policy
\widowpenalty=10000
\clubpenalty=10000
\displaywidowpenalty=10000
\begin{tabularx}{\linewidth}{@{}>{\bfseries}p{0.22\linewidth}X@{}}
組織 & OpenAI / Anthropic \\
種別 & 安全性関連
時系列 & 2025-08-27 (CROSS LAB SAFETY EVALUATION（技術イベント）) \\
\end{tabularx}
\smallskip
{\bfseries 一次資料から整理したtechnical points}
\begin{itemize}[leftmargin=1.5em,itemsep=0.35em]
\item \textbf{一次情報で確認できる事実}: OpenAIとAnthropicは共同評価を公開し、両社が互いのモデルを対象に、ミスアラインメント、指示追従、ハルシネーション、ジェイルブレイクなどの安全性観点を評価した。
\end{itemize}
\begin{minipage}{\linewidth}
% reader-facing Technical Notes generic boundary/source tail group
{\bfseries 読む際の境界}
\begin{itemize}[leftmargin=1.5em,itemsep=0.35em]
\item \textbf{分析上の留意点}: ここで確認できるのは一次資料に記録された範囲の事実であり、ベンダー・プロジェクト・著者による評価を独立再現された結果として扱うものではない。
\end{itemize}
\begin{samepage}
{\bfseries 一次資料}
\begingroup\sloppy
\Urlmuskip=0mu plus 2mu\relax
\begin{itemize}[leftmargin=1.5em,itemsep=0.25em]
\item {\scriptsize\url{https://openai.com/index/openai-anthropic-safety-evaluation}}
\end{itemize}
\endgroup
\end{samepage}
\end{minipage}
\endgroup
\end{technicalnote}
\medskip
